\documentclass[11pt]{article}
\usepackage[paperwidth=210mm,paperheight=297mm,textwidth=182mm,top=20mm,bottom=22mm]{geometry}
\usepackage{amsmath,amssymb,newtxmath}
\usepackage[no-math]{fontspec}
\setmainfont{Times New Roman}
\usepackage{xeCJK}
\setCJKmainfont[AutoFakeBold=2]{SimSun}
\usepackage{graphicx,booktabs,array,hyperref,placeins}
\hypersetup{hidelinks,unicode=true}
\raggedbottom
\setlength{\emergencystretch}{2em}
\renewcommand{\arraystretch}{1.15}
\usepackage{seqsplit}
\begin{document}
\title{Hierarchical Verification, Common-Terminal Current Sharing, and Finite-Length Longitudinal PEEC for Litz-Wire Modeling · 05 Conclusions}
\author{Miao Yuxuan\\Huainan Normal University}
\date{}
\maketitle
\noindent 合并稿对应章节 / Chapter in merged manuscript: \href{paper_en_single.pdf#chapter-05}{05}\par
\FloatBarrier
\section*{05 Conclusions}
\hypertarget{chapter-05}{}
\pdfbookmark[0]{05 Conclusions}{bookmark-chapter-05}
First, the hierarchical framework records shared physical content separately from each model’s omissions, allowing analytical, two-dimensional, longitudinal, and three-dimensional imposed-field results to be used within comparable scopes. Second, the common-terminal augmented system provides a complex-current sharing criterion for the specified linear model, avoiding validation by prescribed equal currents. Third, the executed length and integration sensitivities of finite-length longitudinal PEEC meet the bounded campaign targets, and the final resistance agrees with its equal-area analytical reference. This conclusion remains limited by the longitudinal basis.


The 640-strand two-dimensional field, fixed imposed-field RT0 cuboid, and same-model CPU/GPU comparisons provide frozen evidence that can be redrawn. Three-dimensional electromagnetic ranking of realistic transposed bundles, thermal modeling, and experimental validation remain open. Failure routes retain triggers, original results, and the capabilities required to continue.


\end{document}
